\section[The calculus LJ]{The calculus \(\LJ\)}
\label{sec:lk-lj-04}

The calculus \(\LJ\) uses the same finite-multiset antecedent contexts and the
single-succedent sequent form \(\Gamma\Rightarrow A\).  Its complete primitive
rule surface is as follows.  The zero-premise and left-structural rules are
\[
\frac{}{A\Rightarrow A}\;\textsc{Id}
\qquad
\frac{}{\bottom,\Gamma\Rightarrow C}\;\textsc{\(\bottom L\)},
\]
\[
\frac{\Gamma\Rightarrow C}
     {A,\Gamma\Rightarrow C}\;\textsc{WL}
\qquad
\frac{A,A,\Gamma\Rightarrow C}
     {A,\Gamma\Rightarrow C}\;\textsc{CL}.
\]
Exchange is built into finite-multiset contexts.  The conjunction and
disjunction rules are
\[
\frac{A,B,\Gamma\Rightarrow C}
     {A\land B,\Gamma\Rightarrow C}\;\textsc{\(\land L\)}
\qquad
\frac{\Gamma\Rightarrow A\qquad\Gamma\Rightarrow B}
     {\Gamma\Rightarrow A\land B}\;\textsc{\(\land R\)},
\]
\[
\frac{A,\Gamma\Rightarrow C\qquad B,\Gamma\Rightarrow C}
     {A\lor B,\Gamma\Rightarrow C}\;\textsc{\(\lor L\)}.
\]
The two disjunction-right rules are
\[
\frac{\Gamma\Rightarrow A}
     {\Gamma\Rightarrow A\lor B}\;\textsc{\(\lor R_1\)}
\qquad
\frac{\Gamma\Rightarrow B}
     {\Gamma\Rightarrow A\lor B}\;\textsc{\(\lor R_2\)}.
\]
The implication rules are
\[
\frac{\Gamma\Rightarrow A
      \qquad
      B,\Pi\Rightarrow C}
     {A\to B,\Gamma,\Pi\Rightarrow C}\;\textsc{\(\to L\)}
\qquad
\frac{A,\Gamma\Rightarrow B}
     {\Gamma\Rightarrow A\to B}\;\textsc{\(\to R\)}.
\]
The quantifier rules are
\[
\begin{gathered}
\frac{A[t/x],\Gamma\Rightarrow C}
     {\forall x\,A,\Gamma\Rightarrow C}\;\textsc{\(\forall L\)}
\\[1ex]
\frac{\Gamma\Rightarrow A[y/x]}
     {\Gamma\Rightarrow\forall x\,A}\;\textsc{\(\forall R\)}
\\[1ex]
\frac{A[y/x],\Gamma\Rightarrow C}
     {\exists x\,A,\Gamma\Rightarrow C}\;\textsc{\(\exists L\)}
\\[1ex]
\frac{\Gamma\Rightarrow A[t/x]}
     {\Gamma\Rightarrow\exists x\,A}\;\textsc{\(\exists R\)}.
\end{gathered}
\]
As above, substitution is total and capture avoiding, and \(y\) must not occur
free in the conclusion of \(\forall R\) or \(\exists L\).  There is no right
weakening or right contraction because there is exactly one succedent formula.
Cut is
\[
\frac{\Gamma\Rightarrow A
      \qquad
      A,\Pi\Rightarrow C}
     {\Gamma,\Pi\Rightarrow C}\;\textsc{Cut}.
\]
Negation remains the defined connective
\(\neg A\defeq A\to\bottom\); there are no additional primitive negation
rules.

